\documentclass[11pt,a4paper]{article}
\usepackage{amsmath,amssymb,amsthm}
\usepackage[margin=2.5cm]{geometry}
\usepackage{hyperref}

\newtheorem{definition}{Definition}[section]
\newtheorem{theorem}[definition]{Theorem}
\newtheorem{proposition}[definition]{Proposition}
\newtheorem{remark}[definition]{Remark}
\newtheorem{conjecture}[definition]{Conjecture}

\title{Hyperseed Formalization:\\
  Introducing Hyperseed-1}
\author{ProtomegaTron (automated formalization)}
\date{2026-09-18}

\begin{document}
\maketitle

\begin{abstract}
We formalize the intellectual structure of Ben Goertzel's Substack article
``Introducing Hyperseed-1'' (2024-11-27) within the Hyperseed ontology.
This article introduces the original Hyperseed ontology --- the predecessor
to Hyperseed v2 --- as a semi-formal core ontology designed to seed an
OpenCog Hyperon mind. Key mappings: the seed as initial base space,
semi-formal nature as flexible fiber, inference acceleration as fiber-base
alignment, diverse perspectives as multiple fiber sections, and
hallucinated perspectives as emergent fiber.
\end{abstract}

\tableofcontents

%% =================================================================
\section{The initial base space: seed ontology}
\label{sec:seed}
%% =================================================================

\begin{theorem}[Hyperseed-1 as initial base --- claims 14--18, 27]
\label{thm:seed}
Hyperseed-1 is not trying to encode all commonsense knowledge (the lesson
of Cyc). Instead, it provides a ``seed'' --- a minimal ontological
framework that a reasoning system can use as a starting point for
understanding new concepts.

In Hyperseed terms, Hyperseed-1 is the \emph{initial base space} of the
fiber bundle --- the starting set of semantic primitives that will evolve
as the system learns:
\[
  B(0) = \text{Hyperseed-1 primitives}
\]

The ``seed'' metaphor is literal: it's the initial condition from which the
full base space grows. Just as a biological seed contains the minimal
information needed to grow a complex organism, the ontological seed
contains the minimal conceptual structure needed to bootstrap a reasoning
system.

The seed is designed specifically for an early OpenCog Hyperon mind,
augmented with LLMs or similar technology. It is deliberately semi-formal
--- not fully rigorous predicate logic (too brittle) nor fully informal
natural language (too vague). This semi-formal nature corresponds to a
\emph{flexible fiber}: structured enough to support inference while
flexible enough to evolve:
\[
  \text{Semi-formal} = \text{structured}(\text{inference-support})
  \cap \text{flexible}(\text{evolution-support})
\]

The specific primitive choices include both conventional philosophical
concepts and eccentric ones reflecting Goertzel's particular perspective.
There's a wide range of reasonable options --- the specific list matters
less than the compositional structure.
\end{theorem}

%% =================================================================
\section{Lessons from the history of knowledge formalization}
\label{sec:history}
%% =================================================================

\begin{proposition}[Why not encode everything --- claims 1--5]
\label{prop:history}
The history of commonsense knowledge formalization provides key lessons
that motivate the Hyperseed approach:

\textbf{Leibniz's Universal Characteristic:} The quest to formalize all
human knowledge in logic began with Leibniz and continues to this day.
The dream is a system that captures all of commonsense knowledge in a
formal language.

\textbf{Cyc:} The most substantial practical attempt --- decades of work
encoding commonsense knowledge in predicate logic. Lessons learned:
\begin{enumerate}
\item The number of statements needed is too immense for human encoding.
\item Useful reasoning on massive knowledge bases is computationally hard.
\item Uncertainty management is very difficult when added as afterthought.
\end{enumerate}

\textbf{SUMO:} More elegant and concise than Cyc, but succumbs to the
same fundamental issues.

In Hyperseed terms, these projects attempted to build the \emph{entire
fiber bundle at once} --- encoding all concepts (fibers) over all
situations (base). Hyperseed-1's insight: build only the \emph{base
space} (primitives) and let the fiber grow organically through the
system's experience and reasoning.
\end{proposition}

%% =================================================================
\section{Philosophical foundations: Chalmers and Wierzbicka}
\label{sec:philosophical}
%% =================================================================

\begin{theorem}[Theoretical grounding --- claims 6--13]
\label{thm:foundations}
Two philosophical programs provide theoretical grounding for Hyperseed-1:

\textbf{Chalmers' PQTI framework:} There exists a small set of truths D
(physics P, phenomenal Q, totality T, indexical I) such that a sufficiently
smart reasoner with D plus experience can derive any other truth. This is
not logical derivability --- it requires a reasoning agent to bridge gaps:
\[
  \text{PQTI} + \text{smart reasoner} + \text{experience}
  \vdash \text{any truth}
\]

This inspired by Carnap's Aufbau: begin with a minimal language and a
binary predicate of ``recollected similarity,'' then build everything else.

\textbf{Wierzbicka's NSM:} Approximately 65 semantic primes occur as
single words in every human language. Every concept in any language can be
expressed as a combination of these primes. NSM is empirically grounded
(cross-linguistic analysis of hundreds of languages) but hasn't been
developed into a computational system.

In Hyperseed terms, both Chalmers and Wierzbicka are identifying the
\emph{base space} --- the minimal set of concepts from which all others
can be constructed. Chalmers identifies the base via philosophical
analysis (what's needed for a priori derivability); Wierzbicka identifies
it via empirical linguistics (what's universal across languages).

Hyperseed-1 draws from both but adds the computational dimension: the
base space must not only be theoretically sufficient but must actually
accelerate inference in a computational system.
\end{theorem}

%% =================================================================
\section{Inference acceleration: fiber-base alignment}
\label{sec:acceleration}
%% =================================================================

\begin{proposition}[Ontology as inference accelerator --- claims 19--22, 29]
\label{prop:acceleration}
The practical purpose of Hyperseed-1: speed up inference by providing
conceptual scaffolding. The system doesn't need to derive everything
from scratch.

In Hyperseed terms, the inference speedup comes from \emph{fiber-base
alignment}: when the system encounters a new concept, it can quickly
align it with existing base primitives rather than building understanding
from scratch:
\[
  \text{Time}(\text{understand}(c) | B) \ll
  \text{Time}(\text{understand}(c) | \emptyset)
\]

The ontology provides cross-domain bridges: concepts from different
domains that share ontological primitives are discovered to be
structurally related. A physics concept and an ethics concept that
decompose into similar primitive patterns are structurally analogous
--- and the system discovers this through the base-space alignment.

The ontology is not a taxonomy or classification scheme --- it's meant
to actively drive inference. The distinction: a taxonomy organizes
existing knowledge (passive); an inference accelerator shapes how new
knowledge is acquired and processed (active).
\end{proposition}

%% =================================================================
\section{Multiple perspectives: diverse fiber sections}
\label{sec:perspectives}
%% =================================================================

\begin{proposition}[Diverse and emergent perspectives --- claims 23--26, 30, 32]
\label{prop:perspectives}
The ontology deliberately incorporates multiple philosophical perspectives
rather than committing to one. This includes both established philosophical
viewpoints and ``hallucinated perspectives'' --- perspectives generated by
the ontology construction process itself.

In Hyperseed terms:
\begin{itemize}
\item \textbf{Diverse perspectives = multiple fiber sections:} Different
  philosophical perspectives are different ways of organizing the same base
  material, each providing different inference paths. Including multiple
  sections gives the system multiple routes through concept space.
\item \textbf{Hallucinated perspectives = emergent fiber:} Perspectives
  generated by the construction process are emergent fiber structure ---
  arising from the act of building the bundle, not imported from outside.
  This is creative self-organization: the fiber generates new structure
  through its own construction.
\end{itemize}

The ontology structure consists of:
\begin{enumerate}
\item \textbf{Core entries:} Concepts with formal/informal characterization,
  relationships, and examples.
\item \textbf{Infrastructure concepts:} Meta-concepts about how the ontology
  works (the fiber's self-description).
\item \textbf{Diverse perspective concepts:} Multiple viewpoints on the same
  material (multiple fiber sections).
\end{enumerate}
\end{proposition}

%% =================================================================
\section{Cross-references}
%% =================================================================

\begin{remark}[Connections to other formalizations]
\begin{itemize}
\item \textbf{Hyperseed v2} (2026-03-05): successor. Hyperseed v2 builds on
  Hyperseed-1's intuitions with the fiber bundle formulation, d-calculus,
  and adaptive base space. The informal ``seed'' becomes a formal fiber
  bundle.
\item \textbf{Evidence Is to Logic What Energy Is to Physics} (2026-03-10):
  uncertainty. Cyc's lesson that uncertainty must be baked in connects to
  the evidence conservation framework.
\item \textbf{Tensor Logic} (2025-12-16): computation. Tensor logic provides
  the computational substrate for Hyperseed's fiber operations.
\item \textbf{OpenClaw: Amazing Hands for a Brain} (2026-02-03): brain.
  Hyperseed-1 is a brain improvement (better conceptual structure for
  reasoning), not a hands improvement (better tool interface).
\end{itemize}
\end{remark}

\begin{conjecture}[Seed quality determines growth trajectory]
Conjecture: the quality of the initial seed (base space) determines the
trajectory of the system's conceptual growth --- not its endpoint (which
depends on experience) but the efficiency of the path. A good seed
accelerates growth; a bad seed creates obstacles that must be overcome.

If this conjecture holds, Hyperseed-1 is not just a convenience but a
critical design decision: the initial ontological seed shapes the system's
entire learning trajectory. Different seeds produce systems with different
conceptual ``accents'' --- different ways of carving up the world that
persist even after extensive learning.
\end{conjecture}

\end{document}
